\section{位相の定義}

\subsection{開集合}

\begin{framed}
  \begin{dfn}[位相]
    $S$ を空でない集合とする. $2^S$ の部分集合 $\mathcal{O}$ が以下の条件を満たすとき, $\mathcal{O}$ を $S$ の \textbf{位相} であるといい, 位相 $\mathcal{O}$ が与えられた集合 $S$ を \textbf{位相空間} という.

    \begin{enumerate}
      \item $\displaystyle \emptyset, S\in\mathcal{O}.$
      \item $O_1, O_2\in\mathcal{O}$ ならば $O_1\cap O_2\in\mathcal{O}.$
      \item 任意の添字 $\lambda\in\Lambda$ において $O_\lambda\in\mathcal{O}$ ならば $\displaystyle\bigcup_{\lambda\in\Lambda}O_\lambda\in\mathcal{O}.$
    \end{enumerate}
  \end{dfn}
\end{framed}

位相空間は $(S,\mathcal{O})$ のように書くが, 文脈から与えられている位相が明らかな場合は単に $S$ と書くこともある.
このとき, $S$ を $\textbf{台}$ (または \textbf{台集合}) といい, $S$ の元を \textbf{点} という.

また, $\mathcal{O}$ に属する $S$ の部分集合を \textbf{$\mathcal{O}$-開集合} (あるいは単に \textbf{開集合}) という.

\subsection{閉集合}

$\mathcal{O}$-開集合の補集合を \textbf{$\mathcal{O}$-閉集合} (あるいは単に \textbf{閉集合}) という.


\begin{align}
  \int_{X\times Y}f(x,y)\mu\times\nu(dx,dy)
   & = \int_{X}\left\{\int_Yf(x,y)\nu(dy)\right\}\mu(dx) \notag \\
   & = \int_{Y}\left\{\int_Xf(x,y)\mu(dx)\right\}\nu(dy).
\end{align}